\SourceLeaf{228}%
\[
\cos(1)L=X,\qquad\cos(2)L=Y,\qquad\cos(3)L=Z
\]
\noindent coordinatas puncti \(A\) per \(x,y,z\) denotamus. Sint porro \(x+\mathrm{d}x,y+\mathrm{d}y,z+\mathrm{d}z\) coordinatae alius puncti in superficie curva \(A'\); \(\mathrm{d}s\) ipsius distantia infinite parva ab \(A\); denique \(\lambda\) punctum superficiei sphaericae repraesentans directionem elementi \(AA'\). Erit itaque
\[
\mathrm{d}x=\mathrm{d}s\mathord{.}\cos(1)\lambda,\quad\mathrm{d}y=\mathrm{d}s\mathord{.}\cos(2)\lambda,\quad\mathrm{d}z=\mathrm{d}s\mathord{.}\cos(3)\lambda
\]
\noindent et, quum esse debeat \(\lambda L=90^0\),
\[
X\cos(1)\lambda+Y\cos(2)\lambda+Z\cos(3)\lambda=0
\]
\noindent E combinatione harum aequationum derivamus
\[
X\mathrm{d}x+Y\mathrm{d}y+Z\mathrm{d}z=0
\]
\par\indent Duae habentur methodi generales ad exhibendam indolem superficiei curvae. Methodus \textit{prima} utitur aequatione inter coordinatas \(x,y,z\), quam reductam esse supponemus ad formam \(W=0\), ubi \(W\) erit functio indeterminatarum \(x,y,z\). Sit differentiale completum functionis \(W\)
\[
\mathrm{d}W=P\mathrm{d}x+Q\mathrm{d}y+R\mathrm{d}z
\]
\noindent eritque in superficie curva
\[
P\mathrm{d}x+Q\mathrm{d}y+R\mathrm{d}z=0
\]
\noindent et proin
\[
P\cos(1)\lambda+Q\cos(2)\lambda+R\cos(3)\lambda=0
\]
\noindent Quum haec aequatio, perinde ut ea quam supra stabilivimus, valere debeat pro directionibus omnium elementorum \(\mathrm{d}s\) in superficie curva, facile perspicimus, \(X,Y,Z\) proportionales esse debere ipsis \(P,Q,R\) et proin, quum fiat
\[
XX+YY+ZZ=1
\]
\noindent erit vel
\begingroup\small
\[
X=\frac{P}{\BareRoot{(PP+QQ+RR)}},\quad Y=\frac{Q}{\BareRoot{(PP+QQ+RR)}},\quad Z=\frac{R}{\BareRoot{(PP+QQ+RR)}}
\]\endgroup

\noindent vel
\begingroup\small
\[
X=\frac{-P}{\BareRoot{(PP+QQ+RR)}},\quad Y=\frac{-Q}{\BareRoot{(PP+QQ+RR)}},\quad Z=\frac{-R}{\BareRoot{(PP+QQ+RR)}}
\]\endgroup

